\begin{abstract}
\addchaptertocentry{\tomtatname}
Luận văn đề xuất DT-Guard --- hệ thống phòng thủ chủ động dựa trên Digital Twin cho Federated Learning (FL) trong phát hiện xâm nhập IoT. Khác với các phương pháp phòng thủ thụ động dựa trên phân tích tham số, DT-Guard xây dựng môi trường kiểm chứng hành vi bằng Digital Twin với pipeline kiểm định 4 lớp và Trust-Score. Thông qua A/B Testing năm mô hình sinh dữ liệu, luận văn đánh giá và lựa chọn TabDDPM làm bộ sinh dữ liệu thách thức phù hợp nhất cho môi trường kiểm chứng. Ngoài ra, cơ chế DT-PW (DT-Driven Performance Weighting) với Effort Gate định lượng đóng góp tri thức và loại bỏ free-rider.

Thực nghiệm trên hai bộ dữ liệu CIC-IoT-2023 (34 lớp, 39 đặc trưng) và ToN-IoT (10 lớp, 10 đặc trưng) với 9 phương pháp baseline và 5 loại tấn công (Backdoor, LIE, Min-Max, Min-Sum, MPAF) ở bốn tỉ lệ client độc hại (10\%--50\%) cho thấy: (1) DT-Guard đạt FPR = 0\% trên cả hai tập dữ liệu ở mọi kịch bản; (2) Detection Rate đạt 89,5\%--98,3\% trên CIC-IoT-2023 và 90\%--100\% trên ToN-IoT; (3) DT-PW phát hiện 100\% free-rider (trọng số = 0); (4) Overhead chỉ 80,8 ms/round (\textless 1\%) và peak memory 485 MB. Kết quả cốt lõi đã được đệ trình tại hội nghị IEEE ICCE 2026 {[}23{]}, hiện đang chờ duyệt. Luận văn chứng minh nguyên lý kiểm chứng hành vi chủ động vượt trội hơn phân tích tham số thụ động, và nguyên lý này tổng quát không phụ thuộc dataset hay kiến trúc mô hình IDS cụ thể. Lỗ hổng của các phương pháp thụ động mang tính cấu trúc: cùng một baseline thất bại trước cùng loại tấn công ở cả hai môi trường thử nghiệm.
\end{abstract}
